\documentclass[11pt]{article}
\usepackage{amsmath,amssymb,amsthm}
\usepackage[T1]{fontenc}
\usepackage[utf8]{inputenc}
\usepackage{amsmath,amssymb,amsthm}
\usepackage{enumitem}
\usepackage[margin=1in]{geometry}
\usepackage{hyperref}

\newtheorem{theorem}{Theorem}[section]
\newtheorem{proposition}[theorem]{Proposition}
\newtheorem{lemma}[theorem]{Lemma}
\newtheorem{corollary}[theorem]{Corollary}
\theoremstyle{definition}
\newtheorem{definition}[theorem]{Definition}
\theoremstyle{remark}
\newtheorem{remark}[theorem]{Remark}

\DeclareMathOperator{\Hom}{Hom}
\DeclareMathOperator{\colim}{colim}
\DeclareMathOperator{\Tot}{Tot}
\DeclareMathOperator{\Sh}{Sh}
\DeclareMathOperator{\Ab}{Ab}
\DeclareMathOperator{\cone}{cone}
\newcommand{\gp}{\mathrm{gp}}
\newcommand{\NN}{\mathbb N}
\newcommand{\ZZ}{\mathbb Z}
\newcommand{\RR}{\mathbb R}
\newcommand{\LL}{\mathcal L}
\newcommand{\cG}{\mathcal G}

\newcommand{\fptitle}{Review of Problem 1 -- Astra}
\newcommand{\reviewtag}{Reviewer 1}
\usepackage{fp-review}

\begin{document}
\maketitle

\section{Review}

The proof is correct and very clearly written.  It follows exactly the same strategy as the human solution, but avoids approximation of birhizomic families and therefore achieves a simpler proof.  I ranked it third among the solutions I read.  The similarity to the human solution made this writeup very easy to understand for a reviewer who was familiar with that solution, but this solution might be more difficult to understand for a reader who is less familiar with the problem.  Some basic facts about monoids could have been cited instead of rederived, but since these facts are fairly standard, I do not fault the solution for lack of attributions.

\paragraph{Correctness.}
\begin{itemize}
	\item  1: essentially flawless mathematical argument
	\item  0.5: minor errors that are easily corrected
	\item  0: blank solution, or acknowledged partial result
	\item  $-0.5$: major errors
	\item  $-1$: proposed proof makes no progress towards the result.
\end{itemize}
Note that points should be docked for wrong statements/ solutions.

\noindent Score: 1

\paragraph{Exposition.}
\begin{itemize}
	\item  1: well-written and clear solution.
	\item  0: argument that is possible to follow with serious effort
	\item  $-1$: very poor exposition
\end{itemize}

\noindent Score: 1

\paragraph{Attributions.}
\begin{itemize}
	\item  1: solutions contains fully appropriate attributions
		to related work and theorems being used.
	\item  0: moderate issues with missing attributions or hallucinated references
	\item  $-1$: blatant plagiarism or hallucinations
\end{itemize}

\noindent Score: 1

\newpage
\section{Submission with inline comments}


\begin{abstract}
The answer to the question is affirmative, without a finite-generation
hypothesis on the monoid. We prove a stronger statement: the presheaf
left Kan extension of the given presheaf is already a birhizomic sheaf
and is acyclic on every object of the birhizomic site. The proof uses
finite refinements of birhizomic covers, a finite convex-cover
calculation, and passage to filtered unions of finitely generated
saturated monoids.
\end{abstract}

\section{Statement and conventions}

We use the definition of monoid in the question. In intermediate
constructions, submonoids of groups are allowed to have units; we
explicitly sharpen them before regarding them as objects of the
category in the question. All group-valued expressions involving a
monoid homomorphism refer to its extension to groupifications.

Put
\[
P_x=P[-x]^\sharp\qquad(x\in P).
\]
A presheaf on \(B_P^{\mathrm{op}}\) will be regarded as a covariant
functor \(B_P\to\Ab\). Likewise, presheaves on \(A_P^{\mathrm{op}}\)
are covariant functors on \(A_P\).

For a sheaf \(\mathcal E\) on \(A_P^{\mathrm{op}}\), write
\(H^q(Q,\mathcal E)\) for the right derived functors of evaluation at
the object \(P\to Q\). We make the usual universe choices for sites
and sheaf categories.

\begin{theorem}\label{thm:main}
Let \(P\) be any monoid, and let \(F:B_P\to\Ab\) be any functor.
For \(f:P\to Q\), define
\begin{equation}\label{eq:main-formula}
\cG_F(Q)=
\colim_{\substack{x\in P\\f(x)=0}} F(P_x),
\end{equation}
where the indexing category consists of the corresponding objects
\(P_x\) and their morphisms over \(P\). Then:
\begin{enumerate}[label=\textup{(\roman*)}]
\item \(\cG_F\) is the presheaf left Kan extension of \(F\), and it is
already a sheaf for the birhizomic topology;
\item for every \(P\to Q\),
\[
H^q(Q,\cG_F)=0\qquad(q>0);
\]
\item the canonical unit is a quasi-isomorphism:
\[
F \xrightarrow{\ \sim\ } \mathrm R\pi_*\pi^{-1}F.
\]
\end{enumerate}
\end{theorem}

The indexing category in \eqref{eq:main-formula} is filtered, as will
be proved below. In particular, the formula does not require its
indexing face to be finitely generated.

\section{Elementary constructions with monoids}

\subsection{Filtered colimits and finite rational extensions}

\begin{lemma}\label{lem:filtered}
The groupification of a monoid is torsion-free. Filtered colimits of
monoids are computed on the underlying commutative monoids, and hence
on the underlying sets.
\end{lemma}

\begin{proof}
If \(ng=0\) in \(P^{\gp}\), with \(n>0\), saturation gives
\(g,-g\in P\), and sharpness gives \(g=0\).

Let \(P=\colim_\lambda P_\lambda\) be a filtered colimit of the
underlying commutative monoids. Integrality and sharpness hold in
\(P\): any equality witnessing a failure of either property holds at
some stage, where that property is valid.

To check saturation, write an element of \(P^{\gp}\) as \(a-b\), with
\(a,b\in P\), and suppose \(n(a-b)=c\in P\). By integrality, this is
the equality \(na=nb+c\) in \(P\). Representatives for \(a,b,c\),
together with this equality, exist at a common stage. Saturation at
that stage shows that \(a-b\) belongs to its positive monoid, and
therefore belongs to \(P\).
\end{proof}

\begin{definition}\label{def:rational}
For a finite subset \(S\subset P^{\gp}\), define
\[
C_P(S)=
\left\{g\in P^{\gp}:
ng\in P+\sum_{s\in S}\NN s
\text{ for some }n>0\right\},
\]
and put
\[
P\langle S\rangle=
C_P(S)/\bigl(C_P(S)\cap(-C_P(S))\bigr).
\]
We call \(P\to P\langle S\rangle\) a \emph{finite rational extension}.
This terminology is used only for the construction just specified.
\end{definition}

\begin{lemma}\label{lem:rational}
The construction in Definition~\ref{def:rational} has the following
properties.
\begin{enumerate}[label=\textup{(\roman*)}]
\item For any \(h:P\to M\), there is a morphism
\(P\langle S\rangle\to M\) extending \(h\) if and only if
\(h(s)\in M\) for all \(s\in S\). Such a morphism is unique.
\item The morphism \(P\to P\langle S\rangle\) is of finite
presentation, and its groupification is surjective.
\item Finite rational extensions are preserved by arbitrary
pushout and by composition. Moreover,
\[
P\langle S\rangle\amalg_P P\langle T\rangle
\simeq P\langle S\cup T\rangle.
\]
\item \(P\langle\{-x\}\rangle\simeq P_x\) for \(x\in P\).
\end{enumerate}
\end{lemma}

\begin{proof}
The submonoid \(C=C_P(S)\) is saturated in \(P^{\gp}\), and its units
form the subgroup \(U=C\cap(-C)\). The quotient \(C/U\) is saturated:
if \(n(g+U)\) belongs to \(C/U\), then \(ng\in C\), since adding or
subtracting an element of \(U\) does not change membership in \(C\).
Thus \(g\in C\). The quotient is integral and sharp by construction.

If \(h(s)\in M\) for all \(s\in S\), saturation of \(M\) implies that
\(h(C)\subset M\), and sharpness implies that \(h(U)=0\). This gives
the asserted extension. Its uniqueness follows because
\(P^{\gp}\to(P\langle S\rangle)^{\gp}\) is surjective.

Write \(s=a_s-b_s\), with \(a_s,b_s\in P\). The condition
\(h(s)\in M\) is equivalent to the existence of an element \(z_s\in M\)
such that
\[
h(b_s)+z_s=h(a_s).
\]
There are finitely many such conditions, and each \(z_s\), when it
exists, is unique by integrality. Lemma~\ref{lem:filtered} therefore
shows that
\(\Hom_P(P\langle S\rangle,-)\) commutes with filtered colimits.
This proves finite presentation.

The pushout assertion follows immediately from (i). More generally,
under \(h:P\to Q\), the pushout is
\(Q\langle h(S)\rangle\). For composition, lift the finitely many
group elements used in a rational extension of \(P\langle S\rangle\)
to \(P^{\gp}\); adjoining these lifts together with \(S\) gives the
composite.

Finally, \(h(-x)\in M\), with \(h(x)\in M\), is equivalent to
\(h(x)=0\), since \(M\) is sharp. This is precisely the universal
property of \(P_x\).
\end{proof}

\begin{remark}\label{rem:monomorphisms}
Every finite rational extension is an epimorphism in the category
of monoids, by the uniqueness assertion in
Lemma~\ref{lem:rational}(i). Consequently, its opposite is a
monomorphism of the site. This observation will allow us to use
finite alternating \v{C}ech complexes.
\end{remark}

\subsection{Faces and the category \(B_P\)}

A \emph{face} of \(P\) is a submonoid \(E\subset P\) such that
\(a+b\in E\) implies \(a,b\in E\). The inverse image of \(0\) under a
homomorphism of sharp monoids is a face.

\begin{lemma}\label{lem:faces}
For \(x\in P\), set
\[
E_x=\{a\in P:a\le nx\text{ for some }n\in\NN\}.
\]
Then \(E_x=\ker(P\to P_x)\), and
\[
\Hom_P(P_x,P_y)\ne\varnothing
\quad\Longleftrightarrow\quad E_x\subset E_y.
\]
When this set of morphisms is nonempty, it is a singleton.

If \(P\) is finitely generated, every face is of the form \(E_x\),
there are only finitely many faces, and \(B_P\), up to equivalence,
is their inclusion poset.
\end{lemma}

\begin{proof}
The localization \(P+\NN(-x)\subset P^{\gp}\) is already saturated.
Indeed, if \(ng=a-kx\), with \(a\in P\) and \(k\ge0\), then
\[
n(g+kx)=a+(n-1)kx\in P,
\]
so \(g+kx\in P\). An element \(a\in P\) becomes a unit in this
localization precisely when \(a\le nx\) for some \(n\).
This proves the kernel formula. The assertion about morphisms
follows from the universal property of \(P_x\).

If \(p_1,\ldots,p_r\) generate \(P\), a face \(E\) is generated by
those \(p_j\) which belong to \(E\): every summand in an expression
for an element of a face belongs to that face. Put
\[
x_E=\sum_{p_j\in E}p_j.
\]
Every element of \(E\) is bounded above by a multiple of \(x_E\),
and the face property gives \(E_{x_E}\subset E\). Hence
\(E=E_{x_E}\). The same description proves finiteness of the set of
faces.
\end{proof}

\subsection{Finite-dimensional approximation}

We call a monoid \emph{fs} if it is finitely generated, in addition
to the standing assumptions.

\begin{lemma}\label{lem:fs}
The following assertions hold.
\begin{enumerate}[label=\textup{(\roman*)}]
\item Every monoid is a filtered union of fs submonoids.
\item If \(P\) is fs and \(S\subset P^{\gp}\) is finite, then
\(P\langle S\rangle\) is fs.
\item If \(P\) is fs, there is a homomorphism
\(\ell:P\to\RR_{\ge0}\) which is strictly positive on
\(P\setminus\{0\}\).
\end{enumerate}
\end{lemma}

\begin{proof}
We first recall the relevant elementary lattice calculation.
Let \(L\) be a finitely generated free abelian group, and let
\(v_1,\ldots,v_r\in L\). If
\[
C=\sum_j\RR_{\ge0}v_j,
\]
then the saturation in \(L\) of \(\sum_j\NN v_j\) is \(C\cap L\).
Indeed, a nonnegative real solution to a linear system with rational
coefficients and rational right-hand side can be replaced by a
nonnegative rational solution. One way to see this is to choose a
solution with minimal support: dependence among its supporting
columns permits movement until a coordinate becomes zero; with
independent supporting columns, the solution is rational by linear
algebra. Clearing denominators proves the assertion.

Furthermore, \(C\cap L\) is finitely generated. For
\(g=\sum_j t_jv_j\in C\cap L\), subtract
\(\sum_j\lfloor t_j\rfloor v_j\). The remainder is a lattice point in
the bounded set
\[
\left\{\sum_j u_jv_j:0\le u_j\le1\right\}.
\]
There are only finitely many such remainders, and they, together
with the \(v_j\), generate \(C\cap L\).

For a finite subset of an arbitrary monoid \(P\), take the
saturation of the submonoid it generates in the subgroup it
generates in \(P^{\gp}\). That subgroup is a finitely generated
torsion-free abelian group, hence is free. The saturation lies in
\(P\), is finitely generated by the preceding calculation, and is
sharp. These submonoids are directed under inclusion and exhaust
\(P\). This proves (i). Applying the same calculation to generators
of \(P\) together with \(S\), and then taking the quotient by units,
proves (ii).

For (iii), assume \(P\ne\{0\}\), and choose nonzero generators
\(p_1,\ldots,p_r\). The origin does not belong to their convex hull.
Otherwise the rational-solution argument above would give a
nontrivial relation
\[
\sum_j n_jp_j=0,\qquad n_j\in\NN,
\]
contrary to sharpness. Choose an inner product on
\(P^{\gp}\otimes\RR\), and let \(z\) be the point in this compact
convex hull nearest the origin. Then \(z\ne0\), and minimality gives
\[
\langle z,p_j\rangle\ge \lVert z\rVert^2>0
\]
for every \(j\). The functional \(u\mapsto\langle z,u\rangle\)
has the required property. For \(P=\{0\}\), the assertion is
immediate.
\end{proof}

\section{Infinitesimal envelopes and finite covering refinements}

\subsection{An explicit infinitesimal envelope}

\begin{lemma}\label{lem:envelope}
Let \(v:P\to V\) be a homomorphism to a valuative monoid. There is a
factorization
\[
P\longrightarrow M_v\xrightarrow{p_v}V
\]
such that \(p_v\) is infinitesimal and which is initial among
factorizations of \(v\) whose second morphism is infinitesimal.

More explicitly, with \(G=P^{\gp}\), put
\begin{equation}\label{eq:envelope-cone}
C_v=
\left\{
a\in G:
\begin{array}{l}
\text{there exist }y\in P,\ c\in\NN\text{ such that}\\
a\prec y,\quad v(a)\ge0,\quad v(y)\le c\,v(a)
\end{array}
\right\}.
\end{equation}
Then
\[
M_v=C_v/(C_v\cap(-C_v)).
\]
Moreover,
\begin{equation}\label{eq:envelope-colimit}
M_v=\colim_{\substack{S\subset C_v\\S\text{ finite}}}
P\langle S\rangle
\end{equation}
as monoids under \(P\).
\end{lemma}

\review{This construction was not part of the authors' solution, but was used by the authors in the larger project from which this problem was drawn.}

\begin{proof}
For a nonnegative element \(y\), the condition \(a\prec y\) can
equivalently be written
\[
-ky\le a\le ky
\]
for some \(k\in\NN\). Also, for nonnegative \(b,d\), the condition
\(b\prec d\) is equivalent to \(b\le cd\) for some \(c\in\NN\).

The set \(C_v\) contains \(P\). It is closed under addition: if
\(a,a'\) have witnesses \(y,y'\), then \(a+a'\prec y+y'\), and,
after increasing the two integer bounds to a common bound,
\(v(y+y')\) is bounded by a multiple of \(v(a+a')\).

It is saturated. Suppose \(na\in C_v\), with witness \(y\).
Choose \(k\) such that \(ky\pm na\in P\). Since
\[
n(ky\pm a)=(ky\pm na)+(n-1)ky\in P,
\]
saturation of \(P\) gives \(ky\pm a\in P\).
Saturation of \(V\) gives \(v(a)\ge0\), and the inequality bounding
\(v(y)\) by a multiple of \(nv(a)\) is also a bound by a multiple of
\(v(a)\).

Let \(E=v^{-1}(0)\). We claim that
\begin{equation}\label{eq:envelope-units}
C_v\cap(-C_v)=E^{\gp}.
\end{equation}
If \(a,-a\in C_v\), then \(v(a)=0\). A witness \(y\) for \(a\)
therefore satisfies \(v(y)=0\). For sufficiently large \(k\),
\(a+ky\in P\), and it too has image zero. Thus
\(a=(a+ky)-ky\in E^{\gp}\).
Conversely, a difference of two elements of \(E\) is bounded in
absolute value by their sum and satisfies
\eqref{eq:envelope-cone} in both signs. This proves
\eqref{eq:envelope-units}. Sharpening \(C_v\) gives a monoid \(M_v\)
and a well-defined map \(p_v:M_v\to V\).

We next verify that \(p_v\) is infinitesimal. Let
\(a\in M_v^{\gp}\) and \(b\in M_v\) satisfy
\[
p_v(a)\ge0,\qquad a\prec b,\qquad p_v(b)\prec p_v(a).
\]
Choose representatives \(x\in G\) and \(z\in C_v\) of \(a,b\).
A lower bound in \(a\prec b\) gives an integer \(m\) such that
\[
d=x-mz\in C_v.
\]
Here an equality modulo the unit subgroup can be lifted in this
way because adding or subtracting a unit preserves \(C_v\).

Choose witnesses \(y_d,y_z\in P\) for \(d,z\). The inequality
\(p_v(b)\prec p_v(a)\) gives \(v(z)\le e\,v(x)\) for some \(e\).
Since \(v(d)=v(x)-m v(z)\ge0\), both \(v(d)\) and \(v(z)\) are
bounded by multiples of \(v(x)\). Hence
\(v(y_d+y_z)\) is bounded by a multiple of \(v(x)\).
On the other hand,
\[
x=d+mz\prec y_d+y_z
\]
in the original ordered group \(P^{\gp}\).
Thus \(x\in C_v\), so \(a\in M_v\).

Now consider a factorization
\[
P\xrightarrow{h}H\xrightarrow{p}V
\]
with \(p\) infinitesimal. For \(a\in C_v\), a witness \(y\) in
\eqref{eq:envelope-cone} gives
\[
h(a)\prec h(y),\qquad p(h(a))=v(a)\ge0,\qquad
p(h(y))\prec p(h(a)).
\]
Infinitesimality implies \(h(a)\in H\).
Thus \(h^{\gp}\) maps \(C_v\) into \(H\), kills its units, and
induces \(M_v\to H\). This map is unique because
\(P^{\gp}\to M_v^{\gp}\) is surjective. This proves initiality.

Finally, a homomorphism from the right side of
\eqref{eq:envelope-colimit} to a monoid \(H\) is exactly a
homomorphism \(P\to H\) carrying every element of \(C_v\) into \(H\).
By the universal property of sharpening, this is exactly a
homomorphism \(M_v\to H\). This proves the colimit assertion.
\end{proof}

\begin{corollary}\label{cor:admissible}
If \(S\subset C_v\) is finite, then
\[
P\longrightarrow P\langle S\rangle
\]
is birhizomic at its induced map to \(V\).
\end{corollary}

\begin{proof}
Finite presentation follows from Lemma~\ref{lem:rational}.
For any infinitesimal factorization of \(v\), the proof of
Lemma~\ref{lem:envelope} shows that every element of \(S\) has
nonnegative image in the intermediate monoid. Thus the required
lift exists by Lemma~\ref{lem:rational}(i). It is unique by that
lemma, and its composite to \(V\) is the specified map, again by
uniqueness.
\end{proof}

\subsection{Compactness of valuative tests}

\begin{lemma}\label{lem:valuation-compactness}
Let \(G=P^{\gp}\). The set
\[
K_P=
\{W\subset G:
P\subset W,\ W+W\subset W,\ W\cup(-W)=G\}
\]
is a compact Hausdorff space with the topology induced from
\(\{0,1\}^{G}\).

Every \(W\in K_P\) defines a valuative monoid
\[
V_W=W/(W\cap(-W))
\]
and a homomorphism \(v_W:P\to V_W\). Conversely, every homomorphism
from \(P\) to a valuative monoid defines a point of \(K_P\).
A finite collection of inequalities between group elements defines
a clopen subset of \(K_P\).
\end{lemma}

\begin{proof}
All defining conditions on \(W\) are closed conditions involving
finitely many coordinates at a time. Thus \(K_P\) is closed in the
compact Hausdorff product \(\{0,1\}^G\).

For completeness, compactness of this product follows directly
from the ultrafilter argument. Extend a family of closed subsets
with the finite intersection property to an ultrafilter, using
Zorn's lemma. At each coordinate the ultrafilter chooses exactly
one of the two values. Every basic neighborhood of the resulting
point belongs to the ultrafilter, so that point belongs to every
closed set in the original family.

The submonoid \(W\) is saturated. Indeed, if \(ng\in W\) and
\(g\notin W\), then \(-g\in W\), whence
\[
g=ng+(n-1)(-g)\in W,
\]
a contradiction. Quotienting by its units therefore produces a
valuative monoid. Conversely, for \(v:P\to V\), take
\[
W=\{g\in G:v(g)\ge0\}.
\]
An inequality \(v(a)\ge v(b)\) is the coordinate condition
\(a-b\in W\), proving the last assertion.
\end{proof}

\begin{proposition}[Finite rational refinements]\label{prop:refinement}
Every birhizomic cover of a monoid admits a refinement which is a
finite birhizomic cover by finite rational extensions.

These finite covers form a basis of covering families, stable
under base change and composition.
\end{proposition}

\begin{proof}
Let \(\{P\to Q_i\}_{i\in I}\) be a birhizomic cover.
Fix \(W\in K_P\). Apply the covering condition to
\(v_W:P\to V_W\), and choose
\[
P\xrightarrow{f_i}Q_i\xrightarrow{g_i}V_W
\]
with \(f_i\) birhizomic at \(g_i\).
The infinitesimal map \(M_{v_W}\to V_W\) gives a lift
\[
Q_i\longrightarrow M_{v_W}.
\]
Since \(f_i\) is of finite presentation,
\eqref{eq:envelope-colimit} implies that this lift factors through
\(P\langle S_W\rangle\) for some finite \(S_W\subset C_{v_W}\).
We obtain a refinement morphism
\[
Q_i\longrightarrow P\langle S_W\rangle.
\]

For each \(s\in S_W\), choose a witness \(y_s\in P\), an integer
\(c_s\ge0\), and bounds \(s\prec y_s\), as in
\eqref{eq:envelope-cone}. Let
\[
O_W=
\left\{
W'\in K_P:
s\in W',\ c_s s-y_s\in W'
\text{ for all }s\in S_W
\right\}.
\]
This is a clopen neighborhood of \(W\).
For any valuative homomorphism \(v:P\to V\) whose positive cone
belongs to \(O_W\), the same witnesses show \(S_W\subset C_v\).
Corollary~\ref{cor:admissible} then shows that
\(P\to P\langle S_W\rangle\) is birhizomic at its induced map to
\(V\).

Compactness gives finitely many such neighborhoods covering
\(K_P\). The corresponding rational extensions form a birhizomic
cover, by the preceding observation for arbitrary valuative
targets, and refine the original cover.

Stability under base change and composition follows from
Lemma~\ref{lem:rational} and the topology axioms assumed in the
question.
\end{proof}

\section{The left Kan extension and its transitivity}

\begin{lemma}\label{lem:kan}
For \(F:B_P\to\Ab\), its presheaf left Kan extension to \(A_P\) is
given by \eqref{eq:main-formula}. Its indexing categories are
filtered, and its restriction to \(B_P\) is canonically \(F\).
\end{lemma}

\begin{proof}
By the universal property of \(P_x\), a morphism \(P_x\to Q\) over
\(P\) exists exactly when \(f(x)=0\), and then it is unique.
This gives the usual comma-category formula for left Kan
extension.

If \(x,y\) are killed by \(f\), then \(x+y\) is killed by \(f\), and
there are morphisms
\[
P_x\longrightarrow P_{x+y},
\qquad
P_y\longrightarrow P_{x+y}.
\]
There is also the object \(P_0=P\), and all morphisms in the
indexing category are unique when they exist. Thus the category
is filtered.

For \(Q=P_x\), the comma category has the terminal object
\(P_x\xrightarrow{\mathrm{id}}P_x\), proving the restriction
assertion.
\end{proof}

We use the notation \(\LL_PF\) for this presheaf left Kan extension.

\begin{lemma}[Transitivity]\label{lem:transitivity}
Fix \(f:P\to Q\), and define a presheaf \(H\) on \(B_Q\) by
\[
H(Q_y)=(\LL_PF)(P\to Q\to Q_y).
\]
For every \(Q\to R\), there is a canonical isomorphism
\[
(\LL_QH)(R)\simeq(\LL_PF)(R).
\]
\end{lemma}

\begin{proof}
Expanding the left side gives a filtered colimit over pairs
\((y,x)\) such that
\[
y\in Q,\quad y\mapsto0\text{ in }R,
\qquad
x\in P,\quad x\mapsto0\text{ in }Q_y.
\]
Every such \(x\) is killed in \(R\), so there is a canonical map to
the right side.

Conversely, if \(x\in P\) is killed in \(R\), take \(y=f(x)\).
Then \(y\) is killed in \(R\), and \(x\) is killed in \(Q_y\).
Thus every index on the right occurs on the left.
All finite collections of these indices and their comparisons can
be accommodated at a common larger pair: replace the \(x\)'s and
\(y\)'s by sums, and add the image of the chosen larger \(x\) to
the larger \(y\). Hence any equality in the colimit on the right
already holds at a common index on the left. This proves that the
canonical map is an isomorphism.
\end{proof}

\section{A finite convex-cover calculation}

\subsection{A nerve lemma in the required form}

\review{This reduction to the cohomology of the simplex is the main departure from the human approach.  It is a natural solution to the problem, perhaps more `human' than the human-written version.}

\begin{lemma}\label{lem:convex-nerve}
Let \(U\) be a nonempty convex subset of a finite-dimensional real
vector space. Let \(D_1,\ldots,D_m\) be closed convex polyhedra such
that
\[
U=\bigcup_{i=1}^m(U\cap D_i).
\]
Let \(N\) be the nerve whose nonempty simplices are the subsets
\(I\subset\{1,\ldots,m\}\) for which
\[
U\cap\bigcap_{i\in I}D_i\ne\varnothing.
\]
For every abelian group \(A\), the augmented simplicial cochain
complex of \(N\) with coefficients \(A\) is exact.
\end{lemma}

\review{This is essentially the same simplicial complex used in the authors' solution, but this is presented more concretely.}

\begin{proof}
For each simplex \(I\) of \(N\), choose a point
\(u_I\in U\cap\bigcap_{i\in I}D_i\), and let \(K\) be the convex hull
of these finitely many points. Then \(K\subset U\) is a nonempty
compact convex polytope. The polytopes \(K_i=K\cap D_i\) cover \(K\)
and have exactly the nerve \(N\).

There is a finite triangulation of \(K\) in which every \(K_i\) is
a subcomplex. To construct one, cut \(K\) by the finitely many
hyperplanes defining the \(D_i\), take all resulting polyhedral
cells and their faces, and barycentrically subdivide this finite
polyhedral complex.

For nonempty \(I\), put \(K_I=\bigcap_{i\in I}K_i\). Consider the
double complex
\[
D_{p,q}=
\bigoplus_{i_0<\cdots<i_p}
C_q(K_{\{i_0,\ldots,i_p\}};\ZZ),
\]
with simplicial boundary vertically and the alternating
intersection maps horizontally. The horizontal complex augmented
to \(C_q(K;\ZZ)\) is exact: for each simplex of the triangulation,
its copies are indexed by the nonempty subsets of the nonempty
set of \(i\)'s whose \(K_i\) contain that simplex. This is the
augmented chain complex of a simplex.

Every nonempty \(K_I\) is convex and hence contractible.
Consequently its simplicial homology is \(\ZZ\) in degree zero
and zero in higher degrees. Here one may use homotopy invariance
and the simplicial--singular comparison
\cite[Theorems~2.10 and~2.27]{Hatcher}.
Computing the homology of the finite double complex in its two
directions shows that the nerve has the homology of \(K\).
The double-complex argument is the standard one recalled in
\cite[Tag~012X]{Stacks}.
Thus the augmented integral chain complex of \(N\) is exact.

This is a finite exact complex of finitely generated free abelian
groups, including its augmentation to \(\ZZ\). Its cycle groups
are free, so its short exact sequences split. Applying
\(\Hom_{\ZZ}(-,A)\) therefore preserves exactness and gives the
claimed augmented cochain complex.
\end{proof}

\subsection{Weak valuative covers}

\begin{definition}\label{def:weak}
A finite family
\[
T\longrightarrow T\langle S_i\rangle,\qquad 1\le i\le m,
\]
is \emph{weakly valuative} if every homomorphism from \(T\) to a
valuative monoid factors through at least one member of the
family. No infinitesimal lifting condition is imposed in this
definition.
\end{definition}

Every birhizomic cover is weakly valuative. For a nonempty subset
\(I\subset\{1,\ldots,m\}\), write
\[
T_I=T\left\langle\bigcup_{i\in I}S_i\right\rangle.
\]
These are the intersections in the opposite site.

\begin{proposition}[The fs \v{C}ech calculation]\label{prop:fs-cech}
Suppose \(T\) is fs, \(H:B_T\to\Ab\) is arbitrary, and
\(\{T\to T\langle S_i\rangle\}_{i=1}^m\) is weakly valuative.
Then the augmented alternating complex
\begin{equation}\label{eq:cech}
0\longrightarrow H(T)
\longrightarrow \prod_i(\LL_TH)(T_i)
\longrightarrow \prod_{i<j}(\LL_TH)(T_{\{i,j\}})
\longrightarrow\cdots
\longrightarrow(\LL_TH)(T_{\{1,\ldots,m\}})
\longrightarrow0
\end{equation}
is exact.
\end{proposition}

\review{This is the key finite-type calculation.}

\begin{proof}
By Lemma~\ref{lem:faces}, we may identify \(B_T\) with the finite
poset of faces of \(T\). For a face \(E\), let \(T_E\) denote
\(T_x\) for any \(x\) with \(E_x=E\). If \(R\) is a monoid under
\(T\), its kernel \(E_R=\ker(T\to R)\) is a face, so
\begin{equation}\label{eq:fs-value}
(\LL_TH)(R)=H(T_{E_R}).
\end{equation}

Set \(L=T^{\gp}\), and consider the polyhedral cones
\[
\sigma=\{\ell\in\Hom(L,\RR):\ell(T)\ge0\},
\qquad
\sigma_i=\{\ell\in\sigma:\ell(s)\ge0\text{ for all }s\in S_i\}.
\]
Weak valuative covering, applied to \(\RR_{\ge0}\), implies
\[
\sigma=\bigcup_i\sigma_i.
\]
For nonempty \(I\), put \(\sigma_I=\bigcap_{i\in I}\sigma_i\) and
\(E_I=\ker(T\to T_I)\). The universal property of rational
extensions identifies \(\sigma_I\) with the homomorphisms
\(T_I\to\RR_{\ge0}\). Lemma~\ref{lem:fs} gives
\begin{equation}\label{eq:kernel-detection}
E_I=\{a\in T:\ell(a)=0\text{ for every }\ell\in\sigma_I\}.
\end{equation}
Indeed, a nonzero element of the fs monoid \(T_I\) is detected by a
strictly positive real homomorphism.

Fix a face \(E\) of \(T\) and an abelian group \(A\). Define a
functor \(J_E(A)\) on the face poset by
\[
J_E(A)(T_D)=
\begin{cases}
A,&D\subset E,\\
0,&D\not\subset E,
\end{cases}
\]
with identity transition maps whenever both values are \(A\).

Choose generators \(p_1,\ldots,p_r\) of \(T\), and put
\[
U_E=\{\ell\in\sigma:\ell(p_j)>0
\text{ for every }p_j\notin E\}.
\]
This is nonempty and convex: nonemptiness follows from the
strictly positive functional in Lemma~\ref{lem:fs}, with the case
\(T=\{0\}\) interpreted in the evident way.

We claim that
\begin{equation}\label{eq:nerve-kernel}
E_I\subset E
\quad\Longleftrightarrow\quad
\sigma_I\cap U_E\ne\varnothing.
\end{equation}
For the forward implication, each generator \(p_j\notin E\) is
outside \(E_I\). By \eqref{eq:kernel-detection}, there is an element
of \(\sigma_I\) positive on that generator. Summing these finitely
many functionals gives an element of \(\sigma_I\cap U_E\).
If there are no such generators, the zero functional suffices.

Conversely, if \(E_I\not\subset E\), express an element of
\(E_I\setminus E\) in the chosen generators. Since \(E_I\) is a
face, every generator occurring belongs to \(E_I\); since \(E\)
is a face, at least one occurring generator lies outside \(E\).
Every functional in \(\sigma_I\) vanishes on that generator, so
\(\sigma_I\cap U_E=\varnothing\).

For \(H=J_E(A)\), equations \eqref{eq:fs-value} and
\eqref{eq:nerve-kernel} identify \eqref{eq:cech} with the augmented
cochain complex of the nerve of
\[
U_E=\bigcup_i(U_E\cap\sigma_i).
\]
It is exact by Lemma~\ref{lem:convex-nerve}.

It remains to pass from these functors to arbitrary \(H\).
The functor \(J_E\) is right adjoint to evaluation at \(T_E\):
\[
\Hom(H,J_E(A))=\Hom(H(T_E),A).
\]
Thus \(J_E(A)\) is injective when \(A\) is an injective abelian
group. There are enough such injectives: choose embeddings
\(H(T_E)\hookrightarrow A_E\) into injective abelian groups and
use the adjunctions to obtain
\[
H\hookrightarrow\prod_E J_E(A_E).
\]
The component indexed by \(E=D\) proves injectivity at \(T_D\).
The existence of injective embeddings of abelian groups is
standard; see \cite[Tag~01D6]{Stacks}.
Because the face poset is finite, every injective functor is a
direct summand of a finite product of functors of this form.
Consequently \eqref{eq:cech} is exact for injective \(H\).

Finally, all terms of \eqref{eq:cech}, including its augmentation,
are exact functors of \(H\), by \eqref{eq:fs-value}. Apply them to
an injective resolution of \(H\). The resulting double complex
has exact augmented rows, while its vertical cohomology is the
augmented complex for \(H\). Its horizontal direction is bounded,
and, after shifting the augmented degree by one, it is a
first-quadrant double complex. Computing total cohomology in the
two directions proves exactness for arbitrary \(H\); see again
\cite[Tag~012X]{Stacks}.
\end{proof}

\review{The proof is correct.  It combines several ideas that might better be split into separate statements.}

\section{Passage to arbitrary monoids}

\begin{lemma}\label{lem:weak-descent}
Let \(T=\bigcup_\lambda T_\lambda\) be a filtered union of fs
submonoids. Suppose finite sets \(S_1,\ldots,S_m\subset T^{\gp}\)
define a weakly valuative family over \(T\). After restricting to
a cofinal set of indices, all \(S_i\) lie in \(T_\lambda^{\gp}\),
and
\[
T_\lambda\longrightarrow T_\lambda\langle S_i\rangle,
\qquad 1\le i\le m,
\]
is weakly valuative for every remaining \(\lambda\).
\end{lemma}

\begin{proof}
The finite sets \(S_i\) belong to a common \(T_\lambda^{\gp}\), so
we first restrict to such indices.

The compact spaces of Lemma~\ref{lem:valuation-compactness}
satisfy
\[
K_T\simeq\varprojlim_\lambda K_{T_\lambda}.
\]
Indeed, the groups \(T_\lambda^{\gp}\) have union \(T^{\gp}\),
and compatible positive cones on them give exactly a positive
cone containing \(T\).

In \(K_{T_\lambda}\), put
\[
D_{i,\lambda}=\{W:S_i\subset W\},
\qquad
Z_\lambda=K_{T_\lambda}\setminus\bigcup_iD_{i,\lambda}.
\]
These are clopen sets, and the \(Z_\lambda\) form an inverse
system of compact spaces. If every \(Z_\lambda\) were nonempty,
their inverse limit would be nonempty: every finite collection
of compatibility requirements can be satisfied by choosing a
point at a common upper index, and compactness then supplies a
compatible family. Such a family would define a point of
\(K_T\) lying in none of the \(D_i\), contrary to weak valuative
covering over \(T\).

Thus \(Z_\lambda=\varnothing\) at some stage and at every later
stage. By Lemma~\ref{lem:rational}(i), this is precisely the
desired weak valuative covering condition.
\end{proof}

\begin{proposition}[The general \v{C}ech calculation]\label{prop:general-cech}
Proposition~\ref{prop:fs-cech} holds for an arbitrary monoid \(T\).
\end{proposition}

\review{This is done the same way as in the human solution, but is more explicit.}

\begin{proof}
Choose a filtered union \(T=\bigcup_\lambda T_\lambda\) as in
Lemma~\ref{lem:fs}. By Lemma~\ref{lem:weak-descent}, we may suppose
that the given finite rational family is defined and weakly
valuative at every stage.

For \(x\in T_\lambda\), base change gives
\[
(T_\lambda)_x\amalg_{T_\lambda}T\simeq T_x.
\]
Define \(H_\lambda:B_{T_\lambda}\to\Ab\) by
\[
H_\lambda((T_\lambda)_x)=H(T_x).
\]
For nonempty \(I\subset\{1,\ldots,m\}\), write
\[
T_{I,\lambda}
=T_\lambda\left\langle\bigcup_{i\in I}S_i\right\rangle,
\qquad
T_I=T\left\langle\bigcup_{i\in I}S_i\right\rangle.
\]
The universal property of rational extensions gives
\[
T_I=\colim_\lambda T_{I,\lambda}.
\]

Let \(E_{I,\lambda}=\ker(T_\lambda\to T_{I,\lambda})\), and choose
\(x_{I,\lambda}\in T_\lambda\) generating this face in the sense
of Lemma~\ref{lem:faces}. Then
\[
(\LL_{T_\lambda}H_\lambda)(T_{I,\lambda})
=H(T_{x_{I,\lambda}}).
\]
There is a natural isomorphism
\begin{equation}\label{eq:coefficient-colimit}
(\LL_TH)(T_I)
\simeq
\colim_\lambda
(\LL_{T_\lambda}H_\lambda)(T_{I,\lambda}).
\end{equation}
To verify it, an element \(a\in T\) killed in \(T_I\) belongs to
some \(T_\lambda\) and becomes zero in a later
\(T_{I,\mu}\), by Lemma~\ref{lem:filtered}. It therefore belongs
to \(E_{I,\mu}\), so \(a\le n x_{I,\mu}\) in \(T_\mu\) for some
\(n\). Hence the objects \(T_{x_{I,\lambda}}\) are cofinal in
the indexing category defining \((\LL_TH)(T_I)\).
The same reasoning gives the compatibility maps and proves
\eqref{eq:coefficient-colimit}.

At the augmentation,
\(H_\lambda(T_\lambda)=H(T)\).
There are only finitely many terms in each degree of the
alternating \v{C}ech complex. Therefore
\eqref{eq:coefficient-colimit} identifies the augmented complex
over \(T\) with the filtered colimit of the augmented complexes
over \(T_\lambda\).
The latter are exact by Proposition~\ref{prop:fs-cech}.
Filtered colimits of abelian groups are exact, and they commute
with finite products, so the complex over \(T\) is exact as well;
for exactness of filtered colimits, see
\cite[Tag~00DB]{Stacks}.
\end{proof}

\begin{remark}
Only weak valuative covering was descended in this argument.
There is no need to descend the infinitesimal lifting condition,
and there is no finite-presentation assumption on the coefficient
presheaf.
\end{remark}

\section{From \v{C}ech exactness to acyclicity}

We include the sheaf-theoretic argument to make explicit why the
preceding calculations imply vanishing of derived, rather than
merely \v{C}ech, cohomology.

\begin{lemma}\label{lem:cech-criterion}
Let \(\mathcal C\) be a site admitting a basis of finite covering
families by monomorphisms, stable under base change and
composition, with the corresponding finite intersections.
Suppose an abelian presheaf \(\mathcal E\) has exact augmented
alternating \v{C}ech complex for every covering family in this
basis. Then \(\mathcal E\) is a sheaf and
\[
H^q(U,\mathcal E)=0
\qquad\text{for every }U\in\mathcal C,\ q>0.
\]
\end{lemma}

\begin{proof}
Exactness in augmented degrees \(-1,0\) is the sheaf condition
for the basis. It implies the sheaf condition for all covers:
refine a cover by a basis cover to glue, and use base changes of
that basis cover to verify that the glued section restricts to
the originally specified sections.

We first recall why an injective abelian sheaf \(\mathcal I\)
has exact augmented alternating \v{C}ech complex for a finite
cover by monomorphisms. Let \(\mathbb Z[U]\) denote the free
abelian sheaf on the associated representable sheaf of \(U\).
The augmented chain complex
\[
\cdots\longrightarrow
\bigoplus_{i<j}\mathbb Z[U_i\times_U U_j]
\longrightarrow
\bigoplus_i\mathbb Z[U_i]
\longrightarrow \mathbb Z[U]\longrightarrow0
\]
is exact. Locally over a section of \(U\), one of the covering
subobjects is available, and inserting its index contracts the
augmented complex. This local exactness is exactness after
sheafification. Applying the exact contravariant functor
\(\Hom(-,\mathcal I)\), and using
\(\Hom(\mathbb Z[U],\mathcal I)=\mathcal I(U)\), proves the
assertion. This is also the usual injective \v{C}ech-acyclicity
statement; see \cite[Tag~03AW]{Stacks}.

Embed \(\mathcal E\) in an injective abelian sheaf \(\mathcal I\),
and let \(\mathcal Q=\mathcal I/\mathcal E\).
Such embeddings exist on any site
\cite[Tag~01DP]{Stacks}.
We claim that
\begin{equation}\label{eq:sectionwise-short}
0\longrightarrow\mathcal E(U)\longrightarrow\mathcal I(U)
\longrightarrow\mathcal Q(U)\longrightarrow0
\end{equation}
is exact for every object \(U\).

Indeed, a section of \(\mathcal Q(U)\) lifts locally to
\(\mathcal I\), and we may choose a finite basis cover for these
lifts. Local lifting for epimorphisms of abelian sheaves is part
of the usual description of sheaf cokernels
\cite[Tag~03CM]{Stacks}.
The differences of the local lifts form a \v{C}ech \(1\)-cocycle
with values in \(\mathcal E\). By the assumed exactness, this is a
coboundary. Adjusting the lifts by a \(0\)-cochain in
\(\mathcal E\) makes them agree, and the sheaf condition for
\(\mathcal I\) glues them to a global lift. This proves
\eqref{eq:sectionwise-short}.

For every basis cover, \eqref{eq:sectionwise-short} yields a
short exact sequence of augmented \v{C}ech complexes. The
complexes for \(\mathcal E\) and \(\mathcal I\) are exact, so the
complex for \(\mathcal Q\) is exact. Thus \(\mathcal Q\) has the
same property as \(\mathcal E\).

The long exact sequence of the right derived functors of
evaluation, applied to
\(0\to\mathcal E\to\mathcal I\to\mathcal Q\to0\), gives
\[
H^1(U,\mathcal E)=0,\qquad
H^q(U,\mathcal E)\simeq H^{q-1}(U,\mathcal Q)\quad(q\ge2).
\]
Repeating the argument with \(\mathcal Q\) proves vanishing in
every positive degree.
\end{proof}

\review{This is a nice proof of a standard fact, but could have been omitted.}

\section{Proof of the theorem}

\begin{proof}[Proof of Theorem~\ref{thm:main}]
Let
\[
\cG_F=\LL_PF.
\]
Lemma~\ref{lem:kan} proves formula \eqref{eq:main-formula} and
identifies its restriction to \(B_P\) with \(F\).

Consider any object \(P\to Q\). By
Proposition~\ref{prop:refinement}, every birhizomic cover of this
object has a refinement of the form
\[
Q\longrightarrow Q\langle S_i\rangle,\qquad 1\le i\le m,
\]
which is a finite birhizomic cover. These finite rational covers
form a basis, and their opposites are monomorphisms by
Remark~\ref{rem:monomorphisms}.

Define \(H:B_Q\to\Ab\) by
\[
H(Q_y)=\cG_F(P\to Q\to Q_y).
\]
Lemma~\ref{lem:transitivity} identifies \(\cG_F\), on monoids under
\(Q\), with \(\LL_QH\).
The displayed covering family is weakly valuative. Consequently,
Proposition~\ref{prop:general-cech} proves exactness of its
augmented alternating \v{C}ech complex with coefficients in
\(\cG_F\).

This applies to every object \(Q\) and every basis cover.
Lemma~\ref{lem:cech-criterion} therefore gives both
\[
\cG_F\in\Sh(A_P^{\mathrm{op}},\Ab)
\quad\text{and}\quad
H^q(Q,\cG_F)=0\quad(q>0).
\]
Since \(\cG_F\) is the presheaf left Kan extension and is already
a sheaf, it is canonically \(\pi^{-1}F\). This proves (i) and (ii).

Finally, the right adjoint \(\pi_*\) is restriction to \(B_P\).
Indeed, restriction is right adjoint to presheaf left Kan
extension, and sheafification is left adjoint to inclusion of
sheaves into presheaves. The topology on \(B_P^{\mathrm{op}}\)
is trivial, so no additional sheafification is involved on that
side.

Take an injective resolution \(\cG_F\to\mathcal I^\bullet\) in
abelian sheaves on \(A_P^{\mathrm{op}}\).
Because cohomology in the presheaf category on \(B_P\) is
computed objectwise, for every \(x\in P\) we have
\[
\bigl(\mathrm R^q\pi_*\cG_F\bigr)(P_x)
=
H^q\bigl(\mathcal I^\bullet(P_x)\bigr)
=
H^q(P_x,\cG_F).
\]
These groups vanish for \(q>0\). In degree zero, the canonical
left Kan extension unit gives
\[
(\pi_*\cG_F)(P_x)=\cG_F(P_x)=F(P_x).
\]
Thus the derived unit induces the identity isomorphism in degree
zero and has zero higher cohomology on its target. It is
therefore a quasi-isomorphism:
\[
\boxed{\,F\xrightarrow{\ \sim\ }\mathrm R\pi_*\pi^{-1}F\,}.
\]
\end{proof}

\begin{thebibliography}{9}

\bibitem{Hatcher}
A.~Hatcher,
\emph{Algebraic Topology},
Cambridge University Press, 2002.
Theorems~2.10 and~2.27.

\bibitem{Stacks}
The Stacks Project Authors,
\emph{The Stacks Project}.
Tags \texttt{00DB} (exactness of filtered colimits),
\texttt{012X} (spectral sequences of double complexes),
\texttt{01D6} (injective abelian groups),
\texttt{01DP} (enough injectives for abelian sheaves on a site),
\texttt{03CM} (abelian sheaves and local exactness), and
\texttt{03AW} (injective \v{C}ech acyclicity).

\end{thebibliography}

\end{document}
